\section{Proofs for Path-Space Reduction}
\label{app:path-proofs}

\begin{proof}[Proof of \cref{thm:path-dpp}]
Fix an admissible path \(\gamma\in\Path_{s,t}(x,z)\) and put
\(y=\gamma(r)\).  Restriction and action additivity give
\[
  \mathcal A_{s,t}(\gamma)
  \ge
  c_{s,r}(x,y)+c_{r,t}(y,z).
\]
Taking the infimum over \(\gamma\) proves that the left side of (P1) is at
least the right side.

If the right side is \(+\infty\), the preceding inequality already forces the
left side to be \(+\infty\).  Otherwise, for every \(\varepsilon>0\), choose
\(y\) whose two finite transition costs have sum within \(\varepsilon\) of
the right side, and choose segments \(\gamma_1,\gamma_2\) whose actions are
within \(\varepsilon\) of those costs.  Their concatenation is admissible and
has action at most
\[
  \inf_{w}\{c_{s,r}(x,w)+c_{r,t}(w,z)\}+3\varepsilon.
\]
Let \(\varepsilon\downarrow0\).  The interval lower bounds in
\cref{ass:path-composition} exclude \(-\infty\) transition costs, so every
sum used above is defined.
\end{proof}

\begin{proof}[Proof of \cref{thm:global-mobius-jacobi}]
We divide the argument into the global variational, differential, and Boolean
layers.

\paragraph{Path norm and action regularity.}
For \(\xi\in\mathcal V_u\), parallel transport to one tangent space and the
fundamental theorem of calculus give
\begin{equation}
  \norm{\xi(T)}^2
  \le T\int_0^T\norm{D_t\xi}^2\dd t,
  \qquad
  \int_0^T\norm{\xi(t)}^2\dd t
  \le\frac{T^2}{2}\int_0^T\norm{D_t\xi}^2\dd t.
  \tag{A.P1}
\end{equation}
Hence the derivative norm is equivalent to the usual \(H^1\) norm on the
fixed-initial fluctuation space.  The hypotheses of
\cref{lem:path-action-smoothness} hold with \(k=r+2\); consequently
\((u,\gamma)\mapsto\mathcal A_u(\gamma)\) is a \(C^{r+1}\) functional in
every exponential path chart used below.

\paragraph{Coercivity and existence without a potential lower bound.}
Geodesic convexity gives the global support inequalities
\[
  W_u(t,x)
  \ge W_u(t,x_0)
  +\ip{\operatorname{grad}W_u(t,x_0)}{\log_{x_0}x},
\]
\[
  \Phi_u(x)
  \ge \Phi_u(x_0)
  +\ip{\operatorname{grad}\Phi_u(x_0)}{\log_{x_0}x}.
\]
The uniform base-point bounds imply that both right sides are bounded below
by \(-C(1+d(x_0,x))\), uniformly in \((t,u)\).  Put
\(E(\gamma)=\int_0^T\norm{\dot\gamma}^2\dd t\).  Since
\[
  d(x_0,\gamma(t))\le\sqrt t\,E(\gamma)^{1/2},
\]
there are constants \(C_1,C_2\), independent of \(u\in\mathcal U_0\), such
that
\begin{equation}
  \mathcal A_u(\gamma)
  \ge\frac12E(\gamma)-C_1E(\gamma)^{1/2}-C_2.
  \tag{A.P2}
\end{equation}
Thus minimizing sequences have uniformly bounded \(H^1\) energy even when a
potential tends linearly to \(-\infty\).

The common initial point and the bound
\[
  d(\gamma(s),\gamma(t))
  \le |t-s|^{1/2}E(\gamma)^{1/2}
\]
confine a minimizing sequence to a common bounded ball and make it
equicontinuous.  Finite-dimensional Hadamard manifolds are proper, so
Arzela--Ascoli gives a uniformly convergent subsequence.  In a finite family
of path charts its velocities converge weakly in \(L^2\).  Kinetic energy is
weakly lower semicontinuous, while uniform path convergence and local uniform
continuity of the potentials give convergence of the bulk and terminal
potential terms.  The limit therefore attains the infimum.

\paragraph{Strict path convexity and uniqueness.}
For two admissible paths \(\gamma^0,\gamma^1\), let
\(\Gamma(s,t)\) be their pointwise geodesic homotopy and write
\(H=\partial_s\Gamma\), \(V=\partial_t\Gamma\).  The common initial point
gives \(H(s,0)=0\).  The second variation along the homotopy is
\[
\begin{aligned}
  \frac{\dd^2}{\dd s^2}\mathcal A_u(\Gamma_s)
  ={}&\int_0^T
  \left[\norm{D_tH}^2
  -\ip{R(H,V)V}{H}
  +\operatorname{Hess}W_u[H,H]\right]\dd t\\
  &+\operatorname{Hess}\Phi_u[H(T),H(T)].
\end{aligned}
\]
Every term after \(\norm{D_tH}^2\) is nonnegative by nonpositive curvature
and (MJ2).  If the lower bound vanishes, then \(D_tH=0\); together with
\(H(s,0)=0\), this implies \(H\equiv0\).  Hence the action is strictly convex
along every nontrivial path homotopy.  The minimizer is unique.

\paragraph{Euler--Lagrange equation and index form.}
For \(\xi(0)=0\), the first variation at the minimizer is
\[
\begin{aligned}
  D\mathcal A_u(\gamma_u)[\xi]
  ={}&\int_0^T
  \ip{-D_t\dot\gamma_u+\operatorname{grad}W_u}{\xi}\dd t\\
  &+\ip{\dot\gamma_u(T)+\operatorname{grad}\Phi_u}{\xi(T)}.
\end{aligned}
\]
Compactly supported tests and then arbitrary terminal values give (MJ3) in
weak form.  In a coordinate chart the weak equation first gives
\(\dot\gamma_u\in W^{1,1}\), hence \(\gamma_u\in C^1\); continuity in
\(t\) of the force and smoothness of the connection then give
\(\gamma_u\in C^2\) and (MJ3) pointwise.  Differentiating once more gives
(MJ4).  Nonpositive
curvature and (MJ2) give (MJ5) term by term.  The inequalities in (A.P1) show
that \(\mathfrak B_u\) is continuous and coercive on \(\mathcal V_u\).
Lax--Milgram therefore makes
\(\mathfrak J_u:\mathcal V_u\to\mathcal V_u^*\) an isomorphism, with inverse
norm at most one in the derivative norm.

\paragraph{Smooth global branch.}
In a path chart, write the first-order condition as
\[
  \mathscr F(u,\gamma)=D_\gamma\mathcal A_u(\gamma)=0.
\]
Its path derivative at \(\gamma_u\) is \(\mathfrak J_u\), hence invertible.
The Hilbert implicit-function theorem produces a local \(C^r\) critical
branch around every \(u\in\mathcal U_0\).  Strict convexity makes every
critical path the unique global minimizer, so local branches agree on overlaps
and form one \(C^r\) map on the whole cube neighborhood.

\paragraph{Weak and strong Green--Jacobi response.}
Differentiating \(\mathscr F(u,\gamma_u)=0\) in coordinate \(u_i\) gives
\[
  \mathfrak J_u(\partial_{u_i}\gamma_u)+\mathfrak f_i^u=0,
\]
which proves (MJ6).  For smooth \(\xi,\eta\), covariant integration by parts
in (MJ4) gives
\begin{equation}
\begin{aligned}
  \mathfrak B_u(\xi,\eta)
  ={}&\int_0^T\ip{J_u\xi}{\eta}\dd t\\
  &+\ip{D_t\xi(T)+\operatorname{Hess}\Phi_u\,\xi(T)}{\eta(T)}.
\end{aligned}
\tag{A.P3}
\end{equation}
The closed coercive form corresponds to the Friedrichs self-adjoint operator
with homogeneous conditions
\(\xi(0)=0\) and
\(D_t\xi(T)+\operatorname{Hess}\Phi_u\,\xi(T)=0\).  Comparing
\(\mathfrak B_u(\chi_i^u,\eta)=\mathfrak f_i^u[\eta]\) with (A.P3), first for
compactly supported tests and then for arbitrary endpoint traces, yields the
bulk equation and nonhomogeneous Robin condition in (MJ7).  One-dimensional
elliptic regularity gives the unique classical response.  This also identifies
the weak inverse, Friedrichs resolvent, and Green response.

\paragraph{Value curvature and reciprocity.}
The envelope identity gives
\[
  \partial_{u_i}\mathcal V=\Delta_i(\gamma_u).
\]
Differentiating and using (MJ6) yields
\[
  \partial_{u_i u_j}^2\mathcal V
  =-\mathfrak f_i^u[\chi_j^u]
  =-\mathfrak B_u(\chi_i^u,\chi_j^u).
\]
Symmetry of \(\mathfrak B_u\) proves reciprocity and (MJ8).  Applying the
identity to a linear combination of interventions gives the negative
semidefinite Gram representation.  Under the Riesz identification, the
spectral bounds
\[
  \norm{\mathfrak J_u}^{-1}I
  \preceq\mathfrak J_u^{-1}
  \preceq\norm{\mathfrak J_u^{-1}}I
\]
give (MJ8a).  Since \(\mathfrak J_u^{-1}\) is strictly positive,
\(\langle f,\mathfrak J_u^{-1}f\rangle=0\) holds exactly when \(f=0\).
This proves the kernel and rank identities in (MJ8b).

\paragraph{Boolean integration.}
For a smooth scalar function, the coordinate difference satisfies
\[
  f(u+e_i)-f(u)=\int_0^1\partial_{u_i}f(u+se_i)\dd s
\]
when \(u_i=0\).  Iterating this identity over \(i\in S\) proves (MJ9).
Substitution of (MJ8) proves (MJ10).  Starting the same integration from
\(\mathbf1_C\) proves the conditional-face formula.
\end{proof}

\begin{lemma}[Smoothness of the mechanical action in a path chart]
\label{lem:path-action-smoothness}
Let \(\bar\gamma\in H^1([0,T],\mathcal N)\), and use a smooth exponential
path chart \(\gamma=\exp_{\bar\gamma}\xi\) on an \(H^1\)-neighborhood of
\(\bar\gamma\).  Suppose that the metric coefficients are smooth and that a
bulk potential \(W_u(t,x)\) is continuous in \(t\), \(C^k\) in \((u,x)\),
and has state and parameter derivatives through order \(k\) locally uniform
in \(t\).  Then
\[
  (u,\xi)\longmapsto
  \int_0^T\left(\frac12\norm{\dot\gamma}^2
  +W_u(t,\gamma(t))\right)\dd t
\]
is \(C^{k-1}\) on a sufficiently small path-chart neighborhood.  A
\(C^k\) terminal potential composed with the endpoint trace has the same
regularity.  Its derivatives are the variational derivatives obtained by
differentiating under the integral and at the endpoint.
\end{lemma}

\begin{proof}
The embedding \(H^1([0,T])\hookrightarrow C^0([0,T])\) is continuous, and
one-dimensional \(H^1\) is a Banach algebra.  A small chart ball therefore
maps into one compact tube around \(\bar\gamma([0,T])\).  On that tube, all
coefficient derivatives used below are uniformly bounded.  In local
coordinates the kinetic integrand is
\(g_{ab}(\gamma)\dot\gamma^a\dot\gamma^b/2\).  Repeated differentiation
produces finite sums of bounded Nemytskii coefficients multiplied by at most
two \(L^2\) velocity factors and by \(H^1\) variations.  H\"older's inequality
and the algebra estimate make every derivative through order \(k-1\)
continuous in the \(H^1\) norm.  The potential terms follow from the same
Nemytskii estimate with no velocity factors.  Finally, endpoint evaluation is
a continuous linear map from \(H^1\) to a finite-dimensional tangent space,
so ordinary finite-dimensional composition treats the terminal term.
\end{proof}

\begin{proof}[Proof of \cref{thm:hard-target-jacobi}]
Work in a smooth path chart around \(\gamma_0\) and identify nearby
fixed-initial fluctuation spaces with one Hilbert space \(\mathcal V\).  The
KKT map is
\[
  \mathscr F(u,\gamma,\lambda)
  =\bigl(D_\gamma\mathscr L(u,\gamma,\lambda),
  c_u(\gamma(T))\bigr).
\]
Its derivative in \((\gamma,\lambda)\) at the reference point is precisely
\(\mathfrak K_0\) in \textup{(HT3)}.

We first prove that this derivative is invertible.  Coercivity makes
\(\mathfrak J_0:\mathcal V\to\mathcal V^*\) a positive self-adjoint
isomorphism.  Since \(\mathfrak B_0\ne0\),
\[
  \sigma_0
  =\mathfrak B_0\mathfrak J_0^{-1}\mathfrak B_0^*>0.
\]
For prescribed \((f,a)\in\mathcal V^*\times\R\), the system
\[
  \mathfrak J_0\xi+\mathfrak B_0^*\beta=f,
  \qquad \mathfrak B_0\xi=a
\]
has the unique solution
\[
  \beta=\sigma_0^{-1}
  \bigl(\mathfrak B_0\mathfrak J_0^{-1}f-a\bigr),
  \qquad
  \xi=\mathfrak J_0^{-1}(f-\mathfrak B_0^*\beta).
\]
Hence \(\mathfrak K_0\) is an isomorphism.  The Hilbert implicit-function
theorem gives the unique \(C^r\) KKT branch.  Strict positivity of the
multiplier and coercivity persist after shrinking the parameter
neighborhood.

For a feasible path \(\gamma\) near \(\gamma_u\), Taylor expansion of the
Lagrangian and stationarity give
\[
  \mathcal A_u(\gamma)-\mathcal A_u(\gamma_u)
  =\mathscr L(u,\gamma,\lambda_u)
   -\mathscr L(u,\gamma_u,\lambda_u)
   -\lambda_u c_u(\gamma(T)).
\]
The first difference is positive quadratic to leading order by coercivity,
and the last term is nonnegative because \(\lambda_u>0\) and
\(c_u(\gamma(T))\le0\).  After one uniform neighborhood reduction,
\(\gamma_u\) is therefore the unique constrained local minimizer there.

Differentiating \(\mathscr F(u,\gamma_u,\lambda_u)=0\) gives
\[
  \mathfrak K_u\partial_{u_i}(\gamma_u,\lambda_u)
  +\mathfrak b_i^u=0,
\]
which proves \textup{(HT4)}.  Because the active constraint vanishes,
\[
  \mathcal V_{\rm hard}(u)
  =\mathscr L(u,\gamma_u,\lambda_u).
\]
The envelope identity and a second differentiation yield
\[
  \partial_{u_i u_j}^2\mathcal V_{\rm hard}
  =\partial_{u_i u_j}^2\mathscr L
  +\left\langle
  \mathfrak b_i^u,\partial_{u_j}(\gamma_u,\lambda_u)
  \right\rangle.
\]
Substitution of \textup{(HT4)} proves \textup{(HT5)}.  Self-adjointness of
the bordered map gives reciprocity.  If both data are affine in \(u\), the
direct second derivative is zero; applying the coordinate fundamental theorem
of calculus twice on the indicated face proves \textup{(HT6)}.
\end{proof}

\begin{proof}[Proof of \cref{cor:all-order-mobius-jacobi}]
Apply the set-indexed multivariate Faa di Bruno formula to
\(\mathscr F(u,\gamma_u)=0\).  Since \(\mathcal A_u\), and hence
\(\mathscr F\), is affine in \(u\), at most one index can act on the explicit
\(u\)-argument.  The resulting expansion is
\[
  0=\sum_{\substack{S\subseteq I,\ |S|\le1\\
  \pi\in\Pi(I\setminus S)}}
  D_u^S D_\gamma^{|\pi|}\mathscr F
  \bigl[\eta_B^u\bigr]_{B\in\pi}.
\]
The unique term containing the highest derivative \(\eta_I^u\) corresponds to
\(S=\varnothing\), \(\pi=\{I\}\), and equals
\(\mathfrak J_u\eta_I^u\).  Moving all other terms to the right and applying
\(\mathfrak J_u^{-1}\) proves (MJ11)--(MJ12).  Thus every derivative is
generated recursively from lower-order derivatives by the same Green inverse.

For completeness, applying the same partition formula to the scalar composite
\(\mathcal A_u(\gamma_u)\) gives
\begin{equation}
  \partial_I^{|I|}\mathcal V
  =\sum_{\substack{S\subseteq I,\ |S|\le1\\
  \pi\in\Pi(I\setminus S)}}
  D_u^S D_\gamma^{|\pi|}\mathcal A_u
  \bigl[\eta_B^u\bigr]_{B\in\pi},
  \tag{A.P4}
\end{equation}
where the term \(D_\gamma\mathcal A_u[\eta_I^u]\) vanishes by stationarity.
Equations (A.P4) and (MJ9) give every higher-order Boolean effect.

For three distinct indices, differentiate
\(\mathfrak J_u\partial_{u_j}\gamma_u+D\Delta_j=0\) in \(u_k\), use the
result to eliminate \(\partial_{u_j u_k}^2\gamma_u\) from the second derivative
of \(\partial_{u_i}\mathcal V=\Delta_i(\gamma_u)\), and substitute
\(\partial_{u_l}\gamma_u=-\chi_l\).  Symmetry of \(\mathfrak J_u\) cancels
the second path-response terms and leaves exactly (MJ13).
\end{proof}

\begin{proof}[Proof of \cref{cor:spd-mobius-jacobi}]
The AIRM cone is Hadamard, and a structured realization satisfying the stated
hypotheses has the same intrinsic nonpositive-curvature property.  The
information potentials supply (MJ2), so
\cref{thm:global-mobius-jacobi,cor:all-order-mobius-jacobi} apply directly.
The displayed operator is (MJ7) with the AIRM connection and curvature.  Its
curvature term is nonnegative in the index form, and (MJ10) gives the pair
formula.  The final scope statement records exactly where the classical
Hessian used by the theorem ceases to exist.
\end{proof}

\begin{proof}[Proof of \cref{thm:length-energy-gauge}]
For every finite-energy path, Cauchy--Schwarz gives
\[
  \frac12\int_0^T\norm{\dot\gamma_t}^2\dd t
  \ge
  \frac1{2T}\left(\int_0^T\norm{\dot\gamma_t}\dd t\right)^2
  =\frac{\Len(\gamma)^2}{2T}
  \ge\frac{D(x,\mathcal C)^2}{2T}.
\]
Thus \(\mathcal E_T\ge D^2/(2T)\).  If \(D<+\infty\), choose paths with
length converging to \(D\) and reparameterize each at constant speed.  The
reparameterized paths remain admissible and have energy
\(\Len(\gamma)^2/(2T)\), proving the reverse inequality after taking the
limit.  If \(D=+\infty\), a finite-energy path would have finite length by
Cauchy--Schwarz, a contradiction, so both sides are \(+\infty\).

Equality in Cauchy--Schwarz holds exactly when the metric speed is constant
almost everywhere.  Equality in the second inequality holds exactly for a
length minimizer.  These two equality conditions characterize the energy
minimizers and prove the RDGC specialization.
\end{proof}

\begin{proof}[Proof of \cref{cor:path-bellman}]
Apply \cref{thm:path-dpp} to \(c_{s,T}(x,z)\), add \(\Gamma(z)\), and regroup
the two successive infima:
\[
\begin{aligned}
V(s,x)
&=\inf_z\inf_y
\{c_{s,r}(x,y)+c_{r,T}(y,z)+\Gamma(z)\}\\
&=\inf_y\{c_{s,r}(x,y)+V(r,y)\}.
\end{aligned}
\]
\end{proof}

\begin{proof}[Proof of \cref{thm:path-tonelli}]
Take a minimizing sequence with uniformly bounded objective.  The lower bound
on \(\overline L\) and the lower bound on \(\Gamma\) give a
uniform \(L^p\) bound on path speeds.  Since \(p>1\), H\"older's inequality
gives uniform equicontinuity.  The common initial point and the speed bound
confine every path to a common closed bounded ball.  Finite-dimensional
Hadamard manifolds are proper, so Arzel\`a--Ascoli gives a uniformly convergent
subsequence.  The corresponding velocities converge weakly after passage to
  local charts and a further subsequence.  Sequential viability keeps the limit
  admissible.  More explicitly, subdivide the compact limiting path into
  finitely many intervals whose images lie in normal coordinate charts.  For
  all sufficiently large indices the approximating segments lie in the same
  charts; their coordinate paths converge strongly and uniformly, while their
  coordinate velocities converge weakly in \(L^p\).  On each subinterval the
  normal-integrand hypothesis and convexity in velocity give weak lower
  semicontinuity of the integral functional
  \cite[Chapter~3]{dacorogna2008direct}.  Chart-boundary times form a finite
  measure-zero set, so summing the chartwise inequalities gives the global
  lower bound.  Therefore
\[
  \mathcal A(\gamma)
  \le
  \liminf_n\mathcal A(\gamma_n),
\]
and lower semicontinuity of \(\Gamma\) treats the endpoint.  Hence the limit
attains the value.

If two different minimizing paths existed, their pointwise geodesic
interpolation would be admissible and strict geodesic convexity would give a
strictly smaller objective, a contradiction.
\end{proof}

\begin{proof}[Proof of \cref{prop:path-hjb}]
The Bellman identity on \([t,t+h]\) gives
\[
  V(t,G)=
  \inf_{a_{[t,t+h]}}
  \left\{
  \int_t^{t+h}\ell(s,\gamma_s,a_s)\,\dd s
  +V(t+h,\gamma(t+h))
  \right\}.
\]
  Put
  \[
    q_\varphi(s,X,a)
    =\ell(s,X,a)+\dd_X\varphi(s,X)[b(s,X,a)].
  \]
  Let a smooth test function \(\varphi\) touch \(V\) from above at \((t,G)\).
  For every constant control \(a\in\mathcal U\), the Bellman inequality and the
  local maximum of \(V-\varphi\) give
  \[
    0\le
    \int_t^{t+h}\ell(s,\gamma_s,a)\,\dd s
    +\varphi(t+h,\gamma(t+h))-\varphi(t,G).
  \]
  Uniform local Lipschitz continuity of \(b\) gives
  \(\gamma(t+h)=\exp_G(hb(t,G,a)+o(h))\), uniformly in \(a\).  Divide by
  \(h\), let \(h\downarrow0\), and use compactness of \(\mathcal U\).  The
  inequality holds for every \(a\), hence
  \[
    -\partial_t\varphi(t,G)
    +\mathcal H(t,G,-\dd_G\varphi(t,G))\le0,
  \]
  which is the viscosity subsolution inequality for the sign convention in
  (P4).

  If \(\varphi\) touches \(V\) from below, choose a measurable control whose
  short-horizon cost is within \(o(h)\) of the Bellman infimum.  The local
  minimum of \(V-\varphi\) then yields
  \[
    o(h)\ge
    \int_t^{t+h}\ell(s,\gamma_s,a_s)\,\dd s
    +\varphi(t+h,\gamma(t+h))-\varphi(t,G).
  \]
  Along this trajectory, the chain rule rewrites the right side as
  \[
    \int_t^{t+h}
    \bigl[\partial_s\varphi(s,\gamma_s)
    +q_\varphi(s,\gamma_s,a_s)\bigr]\,\dd s.
  \]
  The trajectory remains uniformly close to \(G\) on the shrinking interval.
  Uniform continuity on the resulting compact neighborhood gives, pointwise in
  the measurable control,
  \[
    q_\varphi(s,\gamma_s,a_s)
    \ge \inf_{a\in\mathcal U}q_\varphi(t,G,a)-o(1).
  \]
  After division by \(h\), this proves
  \[
    -\partial_t\varphi(t,G)
    +\mathcal H(t,G,-\dd_G\varphi(t,G))\ge0.
  \]
  No closure or convexification of averaged velocity--cost pairs is used.
  Continuity of \(V\) and \(\Gamma\) supplies the terminal condition.  At a
  differentiability point, the two inequalities coincide.  The feedback
  formula uses an attaining maximizer, which exists by compactness.  If
  comparison holds in the declared class, standard comparison gives
  uniqueness.
\end{proof}

\begin{proof}[Proof of \cref{thm:path-schur}]
The first-order condition is
\[
  D_\eta\mathfrak A(z,\eta_\star(z))=0.
\]
Coercivity and self-adjointness imply that \(\Jac:\Hilb_0\to\Hilb_0\) is
invertible; equivalently one may apply Lax--Milgram to the second-variation
form.  The Hilbert-space implicit-function theorem therefore produces a unique
\(C^2\) critical section near \((z_0,\eta_0)\).  Coercivity persists after
shrinking the neighborhood, so this section remains the unique local
minimizer.  Differentiating the first-order condition gives
\[
  \mathfrak A_{\eta z}+\Jac D\eta_\star=0,
\]
which proves (P5).

The envelope identity is
\[
  D\Val
  =
  \mathfrak A_z+\mathfrak A_\eta D\eta_\star
  =
  \mathfrak A_z.
\]
Differentiating again and substituting (P5) gives
\[
  D^2\Val
  =
  \mathfrak A_{zz}
  -\mathfrak A_{z\eta}\Jac^{-1}\mathfrak A_{\eta z},
\]
which is (P6).
\end{proof}

\begin{proof}[Proof of \cref{cor:rdgc-jacobi-bridge}]
The classical second-variation formula for kinetic energy along a geodesic
with fixed endpoints is the displayed index form.  Since \(\F\) is totally
geodesic, its intrinsic connection and curvature are the restrictions used by
that formula.  Nonpositive sectional curvature gives
\[
  \ip{R(\dot\gamma_0,h)h}{\dot\gamma_0}\le0,
\]
and hence
\(I(h,h)\ge\int_0^T\norm{D_th}^2\dd t\).  The derivative seminorm is a norm
on \(H_0^1\), equivalent to its standard Sobolev norm by Poincare's inequality.
Thus the Riesz representative of \(I\) is bounded, self-adjoint, and coercive.
The length--energy identification and the final sensitivity claim now follow
from \cref{thm:length-energy-gauge,thm:path-schur}.
\end{proof}

\begin{proof}[Proof of \cref{cor:path-schur-associativity}]
Let \(q(\delta z,\delta\eta_1,\delta\eta_2)\) be the quadratic form defined by
the complete Hessian.  Coercivity gives unique quadratic minimizers.  Hence
\[
  \inf_{\delta\eta_1,\delta\eta_2}q
  =
  \inf_{\delta\eta_1}
  \left(\inf_{\delta\eta_2}q\right).
\]
Each quadratic elimination is represented by the corresponding Schur
complement.  Equality for every \(\delta z\) proves (P7).
\end{proof}

\begin{proof}[Proof of \cref{cor:path-interaction}]
Affineness in \(u\) gives
\[
  \mathfrak A_{uu}=0,
  \qquad
  \mathfrak A_{\eta u}=\mathcal G.
\]
The \(u\)-block of (P6) is therefore
\[
  D^2_{uu}\Val=-\mathcal G^*\Jac^{-1}\mathcal G.
\]
For any \(a\in\R^m\),
\[
  \ip{a}{D^2_{uu}\Val\,a}
  =
  -\ip{\mathcal Ga}{\Jac^{-1}\mathcal Ga}\le0,
\]
which proves (P8) and (P9).  Applying the one-dimensional fundamental theorem
of calculus successively in coordinates \(u_i\) and \(u_j\) gives
\[
  \Val(e_i+e_j)-\Val(e_i)-\Val(e_j)+\Val(0)
  =
  \int_0^1\!\int_0^1
  \partial^2_{u_i u_j}\Val(se_i+te_j)\,\dd s\,\dd t.
\]
Substitution of (P9) proves (P10).
\end{proof}

